\chapter{Разведочный анализ и формирование признаков}

\chaptergoals
  {Показать, как на основе разведочного анализа временного ряда перейти от формального набора наблюдений к осмысленному признаковому пространству, пригодному для дальнейшего прогноза и инженерной интерпретации.}
  {
    \item рассчитывать и интерпретировать описательные статистики;
    \item строить графики, отвечающие на конкретные аналитические вопросы;
    \item выделять календарные, лаговые и агрегированные признаки;
    \item объяснять, почему выбранные признаки содержательно связаны с процессом энергопотребления.
  }

\section{Роль разведочного анализа}

Разведочный анализ необходим не для того, чтобы ``украсить'' исследование несколькими диаграммами, а для того, чтобы понять структуру процесса до начала моделирования. В задачах электроэнергетики это особенно важно, поскольку за таблицей чисел почти всегда стоят реальные режимные закономерности: суточный ход нагрузки, изменение реактивной мощности, влияние календарных факторов, ограниченность диапазонов напряжения и тока.

Если этот этап пропустить, модель будет строиться на формальном наборе полей, а не на содержательно понятном представлении процесса. В результате даже хорошее качество по метрикам может оказаться трудно интерпретируемым с инженерной точки зрения.

\section{Что дают описательные статистики}

Первый шаг разведочного анализа обычно связан с расчетом описательных статистик. Для числовых признаков рассматриваются среднее значение, стандартное отклонение, минимум и максимум. Однако эти величины не следует понимать изолированно от предметного смысла данных.

Например, высокий разброс активной мощности может указывать на выраженную изменчивость нагрузки в течение суток. Узкий диапазон напряжения, наоборот, может говорить о том, что показатель стабилизирован и будет играть в модели иную роль, чем сильно меняющиеся нагрузки. Следовательно, статистика нужна не сама по себе, а как инструмент постановки следующих вопросов.

\section{Почему график должен отвечать на конкретный вопрос}

Хорошая визуализация в учебном блокноте или пособии всегда отвечает на один аналитический вопрос. Для временного ряда такими вопросами могут быть:
\begin{itemize}[leftmargin=1.25cm]
\item как выглядит общий профиль нагрузки во времени;
\item есть ли выраженные пики и провалы;
\item повторяется ли суточный ритм;
\item какие признаки движутся согласованно.
\end{itemize}

Если график не помогает уточнить один из этих вопросов, он остается декоративным. Именно поэтому в связанном практикуме графики дополняются кратким текстом интерпретации: после каждого рисунка объясняется, что именно стало видно и почему это важно для следующего шага.

\screenshotplaceholder
  {fig:eda-load-profile}
  {Пример визуализации временного профиля нагрузки}
  {График временного ряда нагрузки из блокнота \code{03\_eda\_and\_features.ipynb} с читаемыми подписями осей, заголовком и отмеченным ключевым пиком.}
  {Выполнить блок визуализации в \code{03\_eda\_and\_features.ipynb} и сделать скриншот графика так, чтобы были видны подписи осей, заголовок и выделенная характерная точка.}

\figureinstruction
  {Временной профиль нагрузки, используемый для первичной интерпретации энергетического ряда.}
  {После рисунка следует пояснить, что график нужен не для простой демонстрации формы ряда. Он помогает выделить характерные периоды, увидеть максимум нагрузки и подготовить почву для построения календарных и лаговых признаков.}

\section{Переход от наблюдений к признаковому пространству}

После описательных статистик и графиков возникает следующий методический вопрос: какие свойства процесса нужно сохранить в виде признаков для модели. Для временных рядов особенно полезны три группы признаков:
\begin{itemize}[leftmargin=1.25cm]
\item календарные признаки, такие как час суток и день недели;
\item лаговые признаки, связывающие текущее состояние с ближайшей историей;
\item агрегированные признаки, например среднее значение по короткому окну наблюдения.
\end{itemize}

Именно этот переход отличает разведочный анализ от пассивного просмотра графиков. Аналитик не просто описывает данные, а выделяет в них структуру, которая будет передана модели.

\section{Лаговые и агрегированные признаки}

Если обозначить нагрузку как $L_t$, то простейший лаговый признак первого порядка имеет вид
\[
L_{t-1}.
\]
Он показывает, чему равнялось значение в предыдущий момент времени. Для многих энергетических процессов это уже полезная информация, поскольку текущее состояние часто зависит от недавней истории.

Агрегированный признак по окну длины $k$ может быть записан как
\[
\overline{L}_t^{(k)} = \frac{1}{k}\sum_{i=0}^{k-1} L_{t-i}.
\]
Такое усреднение помогает сгладить локальные колебания и выделить более устойчивую структуру ряда.

При построении этих признаков важно не допускать утечки будущей информации. Это означает, что признак, вычисляемый в момент $t$, должен опираться только на текущие и прошлые наблюдения, но не на будущие значения ряда.

\section{Типичные ошибки при формировании признаков}

Наиболее частые ошибки на этом этапе связаны не с синтаксисом, а с логикой:
\begin{itemize}[leftmargin=1.25cm]
\item использование признаков, смысл которых не объяснен предметно;
\item механическое добавление большого числа полей без проверки их полезности;
\item утечка будущей информации при построении лагов и скользящих статистик;
\item попытка интерпретировать корреляцию как причинную связь.
\end{itemize}

Методически правильный подход состоит в том, чтобы каждый признак можно было объяснить двумя способами: математически и инженерно. Если признак трудно описать на языке задачи, его роль в модели требует дополнительного обоснования.

\section{Связь главы с практикумом}

Главе соответствует блокнот \notebookhref{03_eda_and_features.ipynb}{\code{03\_eda\_and\_features.ipynb}}. В нем последовательно показаны диапазон наблюдений, описательные статистики, график нагрузки, суточный профиль и корреляционная структура признаков. За счет этого читатель видит, как из исходного временного ряда постепенно формируется пространство признаков для следующего этапа анализа.
